\documentclass[10pt,a4paper]{amsart}
\usepackage{amsmath,amssymb,amsthm,graphicx}
% Retained minimal preamble for the Stacks Project chapter varieties.tex.
% Verbatim excerpt of the chapter preamble: Xy-pic setup (native preamble lines 14-18), the native theorem
% declarations with the mathematical macros (native preamble lines 154-217), and the project own no-op
% definitions of the reference/slogan/history display environments (native wrapper.tex lines 1-3) that the
% counted statement uses for its slogan. Copied mechanically; nothing added, removed or redefined.
% For commutative diagrams we use Xy-pic
\usepackage[all]{xy}

% We use 2cell for 2-commutative diagrams.
\xyoption{2cell}
% Theorem environments.
%
\theoremstyle{plain}
\newtheorem{theorem}[subsection]{Theorem}
\newtheorem{proposition}[subsection]{Proposition}
\newtheorem{lemma}[subsection]{Lemma}

\theoremstyle{definition}
\newtheorem{definition}[subsection]{Definition}
\newtheorem{example}[subsection]{Example}
\newtheorem{exercise}[subsection]{Exercise}
\newtheorem{situation}[subsection]{Situation}

\theoremstyle{remark}
\newtheorem{remark}[subsection]{Remark}
\newtheorem{remarks}[subsection]{Remarks}

\numberwithin{equation}{subsection}

% Macros
%
\def\lim{\mathop{\mathrm{lim}}\nolimits}
\def\colim{\mathop{\mathrm{colim}}\nolimits}
\def\Spec{\mathop{\mathrm{Spec}}}
\def\Hom{\mathop{\mathrm{Hom}}\nolimits}
\def\Ext{\mathop{\mathrm{Ext}}\nolimits}
\def\SheafHom{\mathop{\mathcal{H}\!\mathit{om}}\nolimits}
\def\SheafExt{\mathop{\mathcal{E}\!\mathit{xt}}\nolimits}
\def\Sch{\mathit{Sch}}
\def\Mor{\mathop{\mathrm{Mor}}\nolimits}
\def\Ob{\mathop{\mathrm{Ob}}\nolimits}
\def\Sh{\mathop{\mathit{Sh}}\nolimits}
\def\NL{\mathop{N\!L}\nolimits}
\def\CH{\mathop{\mathrm{CH}}\nolimits}
\def\proetale{{pro\text{-}\acute{e}tale}}
\def\etale{{\acute{e}tale}}
\def\QCoh{\mathit{QCoh}}
\def\Ker{\mathop{\mathrm{Ker}}}
\def\Im{\mathop{\mathrm{Im}}}
\def\Coker{\mathop{\mathrm{Coker}}}
\def\Coim{\mathop{\mathrm{Coim}}}

% Boxtimes
%
\DeclareMathSymbol{\boxtimes}{\mathbin}{AMSa}{"02}

%
% Macros for moduli stacks/spaces
%
\def\QCohstack{\mathcal{QC}\!\mathit{oh}}
\def\Cohstack{\mathcal{C}\!\mathit{oh}}
\def\Spacesstack{\mathcal{S}\!\mathit{paces}}
\def\Quotfunctor{\mathrm{Quot}}
\def\Hilbfunctor{\mathrm{Hilb}}
\def\Curvesstack{\mathcal{C}\!\mathit{urves}}
\def\Polarizedstack{\mathcal{P}\!\mathit{olarized}}
\def\Complexesstack{\mathcal{C}\!\mathit{omplexes}}
% \Pic is the operator that assigns to X its picard group, usage \Pic(X)
% \Picardstack_{X/B} denotes the Picard stack of X over B
% \Picardfunctor_{X/B} denotes the Picard functor of X over B
\def\Pic{\mathop{\mathrm{Pic}}\nolimits}
\def\Picardstack{\mathcal{P}\!\mathit{ic}}
\def\Picardfunctor{\mathrm{Pic}}
\def\Deformationcategory{\mathcal{D}\!\mathit{ef}}
\newenvironment{reference}{}{}
\newenvironment{slogan}{}{}
\newenvironment{history}{}{}
\usepackage{cleveref}
\crefname{originallemma}{original lemma}{original lemmas}
\Crefname{originallemma}{Original lemma}{Original lemmas}
\makeatletter
\newlabel{lemma-Galois-action-quasi-compact-open}{{04KU}{source}{}{sourcepointer-1}{}}
\newlabel{lemma-Galois-action-quasi-compact-open@cref}{{[originallemma][0][]04KU}{[1][0][]source}{}{}{}}
\makeatother
\crefname{originallemma}{original lemma}{original lemmas}
\Crefname{originallemma}{Original lemma}{Original lemmas}
\makeatletter
\newlabel{lemma-characterize-geometrically-connected}{{0387}{source}{}{sourcepointer-2}{}}
\newlabel{lemma-characterize-geometrically-connected@cref}{{[originallemma][0][]0387}{[1][0][]source}{}{}{}}
\makeatother
\crefname{originallemma}{original lemma}{original lemmas}
\Crefname{originallemma}{Original lemma}{Original lemmas}
\makeatletter
\newlabel{lemma-closed-fixed-by-Galois}{{038B}{source}{}{sourcepointer-3}{}}
\newlabel{lemma-closed-fixed-by-Galois@cref}{{[originallemma][0][]038B}{[1][0][]source}{}{}{}}
\makeatother
\title{Characterization of geometric connectedness by connectedness after every finite separable field extension}
\author{The Stacks Project Authors}
\date{stacks-project:varieties.tex:local-snapshot-20261011; GNU FDL 1.2 or later}
\begin{document}
\maketitle
\noindent\textit{Editorial note.} All delivered mathematical text is the author's own, copied from the pinned Stacks Project chapter file varieties.tex; no statement or proof step was rewritten, reordered or removed.\par\medskip
\noindent\textit{Editorial note.} The proof contains the sentence 'Verification omitted.' at the compatibility claim for quasi-compact opens W contained in W'; the sentence is reproduced verbatim. The omitted check is a bounded compatibility statement and the surrounding descent construction is complete as written.\par\medskip
\noindent\textit{Editorial note.} The counted spans carry 3 navigation stubs and no bibliography citation.\par\medskip
\noindent\textit{Editorial note.} The preamble is the chapter's retained minimal preamble (sources/stacks-varieties/retained-preamble.tex lines 1-77), a verbatim excerpt of the native preamble.\par\medskip
\noindent\textit{Editorial note.} The label lemma-characterize-geometrically-disconnected is defined only in varieties.tex, so no cross-chapter restatement had to be weighed for this card.\par\medskip
\noindent\textit{Allowed standard prerequisites (editorial).}\begin{itemize}
\item Geometrically connected scheme. Same chapter, Lemma lemma-characterize-geometrically-connected, invoked at proof line 952. Supplies that failure of geometric connectedness gives a disconnected X over the separable closure.
\item Galois action on quasi-compact opens of the base change to the separable closure. Same chapter, Lemma lemma-Galois-action-quasi-compact-open, invoked at proof lines 962-967. Supplies that the pieces of a quasi-compact open are defined over a finite subextension and fixed by an open subgroup.
\item Closed subsets of the base change that are fixed by the Galois action descend. Same chapter, Lemma lemma-closed-fixed-by-Galois, invoked at proof lines 1012-1014. Supplies that the Galois invariant pieces U and V come from closed subschemes over a finite subextension.
\end{itemize}\medskip
\ifdefined\hypertarget\hypertarget{sourcepointer-1}{}\fi
\noindent\textit{Source reference.} Original lemma 04KU: sources/stacks-varieties/pinned-original-source.tex, source lines 877--877. Pointer outside the counted span, invoked once for context; the load-bearing argument of the counted statement is the author's proof reproduced in full in the proof span.\par\medskip
\ifdefined\hypertarget\hypertarget{sourcepointer-2}{}\fi
\noindent\textit{Source reference.} Original lemma 0387: sources/stacks-varieties/pinned-original-source.tex, source lines 792--792. Pointer outside the counted span, invoked once for context; the load-bearing argument of the counted statement is the author's proof reproduced in full in the proof span.\par\medskip
\ifdefined\hypertarget\hypertarget{sourcepointer-3}{}\fi
\noindent\textit{Source reference.} Original lemma 038B: sources/stacks-varieties/pinned-original-source.tex, source lines 899--899. Pointer outside the counted span, invoked once for context; the load-bearing argument of the counted statement is the author's proof reproduced in full in the proof span.\par\medskip
\begin{lemma}
\label{lemma-characterize-geometrically-disconnected}
Let $k$ be a field.
Let $X$ be a scheme over $k$.
The following are equivalent
\begin{enumerate}
\item $X$ is geometrically connected,
\item for every finite separable field extension $k'/k$
the scheme $X_{k'}$ is connected.
\end{enumerate}
\end{lemma}

\begin{proof}
It follows immediately from the definition that (1) implies (2).
Assume that $X$ is not geometrically connected.
Let $k \subset \overline{k}$ be a separable algebraic
closure of $k$. By
Lemma \ref{lemma-characterize-geometrically-connected}
it follows that $X_{\overline{k}}$ is disconnected.
Say $X_{\overline{k}} = \overline{U} \amalg \overline{V}$
with $\overline{U}$ and $\overline{V}$ open, closed, and nonempty.

\medskip\noindent
Suppose that $W \subset X$ is any quasi-compact open.
Then $W_{\overline{k}} \cap \overline{U}$ and
$W_{\overline{k}} \cap \overline{V}$ are open and closed in
$W_{\overline{k}}$. In particular $W_{\overline{k}} \cap \overline{U}$ and
$W_{\overline{k}} \cap \overline{V}$ are quasi-compact, and by
Lemma \ref{lemma-Galois-action-quasi-compact-open}
both $W_{\overline{k}} \cap \overline{U}$ and
$W_{\overline{k}} \cap \overline{V}$
are defined over a finite subextension and invariant under an
open subgroup of $\text{Gal}(\overline{k}/k)$.
We will use this without further mention in the following.

\medskip\noindent
Pick $W_0 \subset X$ quasi-compact open such that both
$W_{0, \overline{k}} \cap \overline{U}$ and
$W_{0, \overline{k}} \cap \overline{V}$ are nonempty.
Choose a finite subextension $\overline{k}/k'/k$
and a decomposition $W_{0, k'} = U_0' \amalg V_0'$ into open and closed
subsets such that
$W_{0, \overline{k}} \cap \overline{U} = (U'_0)_{\overline{k}}$ and
$W_{0, \overline{k}} \cap \overline{V} = (V'_0)_{\overline{k}}$.
Let $H = \text{Gal}(\overline{k}/k') \subset \text{Gal}(\overline{k}/k)$.
In particular
$\sigma(W_{0, \overline{k}} \cap \overline{U}) =
W_{0, \overline{k}} \cap \overline{U}$ and similarly for
$\overline{V}$.

\medskip\noindent
Having chosen $W_0$, $k'$ as above, for every quasi-compact open
$W \subset X$ we set
$$
U_W =
\bigcap\nolimits_{\sigma \in H} \sigma(W_{\overline{k}} \cap \overline{U}),
\quad
V_W =
\bigcup\nolimits_{\sigma \in H} \sigma(W_{\overline{k}} \cap \overline{V}).
$$
Now, since $W_{\overline{k}} \cap \overline{U}$ and
$W_{\overline{k}} \cap \overline{V}$ are fixed by an open subgroup of
$\text{Gal}(\overline{k}/k)$ we see that the union and intersection
above are finite. Hence $U_W$ and $V_W$ are both open and closed.
Also, by construction $W_{\bar k} = U_W \amalg V_W$.

\medskip\noindent
We claim that if $W \subset W' \subset X$ are quasi-compact
open, then $W_{\overline{k}} \cap U_{W'} = U_W$ and
$W_{\overline{k}} \cap V_{W'} = V_W$. Verification omitted.
Hence we see that upon defining $U = \bigcup_{W \subset X} U_W$
and $V = \bigcup_{W \subset X} V_W$ we obtain
$X_{\overline{k}} = U \amalg V$ is a disjoint union of open
and closed subsets.
It is clear that $V$ is nonempty as it is constructed by taking
unions (locally). On the other hand, $U$ is nonempty since it contains
$W_0 \cap \overline{U}$ by construction. Finally, $U, V \subset X_{\bar k}$
are closed and $H$-invariant by construction. Hence by
Lemma \ref{lemma-closed-fixed-by-Galois}
we have $U = (U')_{\bar k}$, and $V = (V')_{\bar k}$ for some
closed $U', V' \subset X_{k'}$. Clearly $X_{k'} = U' \amalg V'$
and we see that $X_{k'}$ is disconnected as desired.
\end{proof}

\end{document}
